% GENERATED FILE. Edit canonical lessons, then run npm run generate:books.
% canonical-source-hash: fnv1a64:80a3ba5378b192ba
% canonical-lessons: BN-C50-dhoya, BN-C50-hasa, BN-C50-kada, BN-C50-gaowa, BN-C50-nacha

\chapter{To wash, To laugh, To cry, To sing, To dance}
\label{ch:bn-to-wash-to-laugh-to-cry-to-sing-to-dance}
\hlchaptermodality{50}

\begin{chapteropening}
\textbf{By the end of this chapter:} \emph{I can say to wash, to laugh, to cry, to sing and to dance.}

Complete the last lesson of chapter 50: i can say to wash, to laugh, to cry, to sing and to dance.
\end{chapteropening}

\section[dhoyā]{\bn{ধোয়া} (dhoyā) — to wash}

\label{lesson:BN-C50-dhoya}



\begin{quote}
Before the new one: say the Bengali for to put, then the Bengali for to open.
\end{quote}

\subsection*{You'll want to know: \bn{ধোয়া}}
\textbf{\bn{ধোয়া}} — \emph{dhoyā} — \textquotedblleft{}to wash\textquotedblright{}.

\begin{grammarlens}[title={What you've built}]
One more word for this chapter.
\end{grammarlens}

\subsection*{Your turn}
\begin{itemize}
\raggedright
  \item \emph{Say it:} \emph{dhoyā}
  \item \emph{Say it:} \emph{dhoyā}, once more
  \item \emph{Say it:} say \emph{dhoyā}
  \item \emph{Recall:} say the Bengali for to put, then the Bengali for to open, then say \emph{dhoyā} again
\end{itemize}

\begin{tcolorbox}[breakable,colback=teal!4,colframe=teal!35!black,title={Before you move on}]
What does \bn{ধোয়া} mean? (\textquotedblleft{}To wash\textquotedblright{}.) Say it once more.
\end{tcolorbox}
\section[hāsā]{\bn{হাসা} (hāsā) — to laugh}

\label{lesson:BN-C50-hasa}



\begin{quote}
Before the new one: say the Bengali for to open, then the Bengali for to wash.
\end{quote}

\subsection*{You'll want to know: \bn{হাসা}}
\textbf{\bn{হাসা}} — \emph{hāsā} — \textquotedblleft{}to laugh\textquotedblright{}.

\begin{grammarlens}[title={What you've built}]
One more word for this chapter.
\end{grammarlens}

\subsection*{Your turn}
\begin{itemize}
\raggedright
  \item \emph{Say it:} \emph{hāsā}
  \item \emph{Say it:} \emph{hāsā}, once more
  \item \emph{Say it:} say \emph{hāsā}
  \item \emph{Recall:} say the Bengali for to open, then the Bengali for to wash, then say \emph{hāsā} again
\end{itemize}

\begin{tcolorbox}[breakable,colback=teal!4,colframe=teal!35!black,title={Before you move on}]
What does \bn{হাসা} mean? (\textquotedblleft{}To laugh\textquotedblright{}.) Say it once more.
\end{tcolorbox}
\section[kā̃dā]{\bn{কাঁদা} (kā̃dā) — to cry}

\label{lesson:BN-C50-kada}



\begin{quote}
Before the new one: say the Bengali for to wash, then the Bengali for to laugh.
\end{quote}

\subsection*{You'll want to know: \bn{কাঁদা}}
\textbf{\bn{কাঁদা}} — \emph{kā̃dā} — \textquotedblleft{}to cry\textquotedblright{}.

\begin{grammarlens}[title={What you've built}]
One more word for this chapter.
\end{grammarlens}

\subsection*{Your turn}
\begin{itemize}
\raggedright
  \item \emph{Say it:} \emph{kā̃dā}
  \item \emph{Say it:} \emph{kā̃dā}, once more
  \item \emph{Say it:} say \emph{kā̃dā}
  \item \emph{Recall:} say the Bengali for to wash, then the Bengali for to laugh, then say \emph{kā̃dā} again
\end{itemize}

\begin{tcolorbox}[breakable,colback=teal!4,colframe=teal!35!black,title={Before you move on}]
What does \bn{কাঁদা} mean? (\textquotedblleft{}To cry\textquotedblright{}.) Say it once more.
\end{tcolorbox}
\section[gāowā]{\bn{গাওয়া} (gāowā) — to sing}

\label{lesson:BN-C50-gaowa}



\begin{quote}
Before the new one: say the Bengali for to laugh, then the Bengali for to cry.
\end{quote}

\subsection*{You'll want to know: \bn{গাওয়া}}
\textbf{\bn{গাওয়া}} — \emph{gāowā} — \textquotedblleft{}to sing\textquotedblright{}.

\begin{grammarlens}[title={What you've built}]
One more word for this chapter.
\end{grammarlens}

\subsection*{Your turn}
\begin{itemize}
\raggedright
  \item \emph{Say it:} \emph{gāowā}
  \item \emph{Say it:} \emph{gāowā}, once more
  \item \emph{Say it:} say \emph{gāowā}
  \item \emph{Recall:} say the Bengali for to laugh, then the Bengali for to cry, then say \emph{gāowā} again
\end{itemize}

\begin{tcolorbox}[breakable,colback=teal!4,colframe=teal!35!black,title={Before you move on}]
What does \bn{গাওয়া} mean? (\textquotedblleft{}To sing\textquotedblright{}.) Say it once more.
\end{tcolorbox}
\section[nāchā]{\bn{নাচা} (nāchā) — to dance}

\label{lesson:BN-C50-nacha}



\begin{quote}
Before the new one: say the Bengali for to cry, then the Bengali for to sing.
\end{quote}

\subsection*{You'll want to know: \bn{নাচা}}
\textbf{\bn{নাচা}} — \emph{nāchā} — \textquotedblleft{}to dance\textquotedblright{}.

\begin{grammarlens}[title={What you've built}]
One more word for this chapter.
\end{grammarlens}

\subsection*{Your turn}
\begin{itemize}
\raggedright
  \item \emph{Say it:} \emph{nāchā}
  \item \emph{Say it:} \emph{nāchā}, once more
  \item \emph{Say it:} say \emph{nāchā}
  \item \emph{Recall:} say the Bengali for to cry, then the Bengali for to sing, then say \emph{nāchā} again
\end{itemize}

\begin{tcolorbox}[breakable,colback=teal!4,colframe=teal!35!black,title={Before you move on}]
What does \bn{নাচা} mean? (\textquotedblleft{}To dance\textquotedblright{}.) Say it once more.
\end{tcolorbox}
